\section{서 론}

기존 온라인 데이팅 서비스에서 사용자 매칭은 주로 나이, 거주 지역, 관심사나 간단한 자기소개 문구 같은 텍스트 프로필 정보에 의존한다. 그러나 실제 사용자는 패션 스타일, 사진 분위기, 촬영 환경과 같은 시각적 요인에 크게 좌우되며 단순 분류 모델(Classification)이나 필터링으로는 이러한 고차원적인 취향을 충분히 반영하기 어렵다.

이러한 한계를 극복하기 위해 비전-언어 모델(Vision-Language Model) 기반의 특징 추출과 Metric Learning을 결합한 스타일 매칭 시스템을 제안한다. 이를 통해 텍스트 필터로 포착할 수 없는 `취향'과 `분위기'를 이미지 임베딩 벡터의 코사인 유사도(Cosine Similarity)를 기반으로 정량화하여 해결한다.

본 연구는 Walster 등이 제안한 사회심리학적 이론인 `매칭 가설(The Matching Hypothesis)\cite{walster1966}'을 이론적 토대로 한다. 해당 가설에 따르면, 개인은 파트너 선택 시 자신과 유사한 수준의 신체적 매력도나 사회적 바람직성을 가진 상대를 선호하는 경향(Assortative Mating)이 있다. 이를 현대적 온라인 데이팅 환경에 적용하여, 텍스트 기반 필터링으로는 포착하기 어려운 시각적 스타일과 분위기의 유사성을 매칭의 핵심 요소로 간주하고, 임베딩 공간상의 거리로 정량화함으로써 사용자 간의 매칭 만족도를 향상시키는 것을 목표로 한다.
